\section{Evaluation}
\label{sec:evaluation}

We ask five questions. Is DriftGate more accurate than existing inference rules under mobility (Section~\ref{sec:overall})? Which devices, requests, and times of day does the gain come from (Section~\ref{sec:where})? How much accuracy does it keep when fewer requests are offloaded (Section~\ref{sec:eval_offload})? Which parts of the design produce the gain (Section~\ref{sec:ablation})? And how does it behave at larger scale, for the least accurate devices, and on real devices (Sections~\ref{sec:scale} and~\ref{sec:device})?

\subsection{Setup}
\label{sec:setup}

\textbf{Settings.} Table~\ref{tab:settings} lists the settings. The commute scenario S1 places five cells with a 1\,km edge range, four residential cells and one hub, and 50 devices whose users walk from home to the hub and back over a day from 05:00 to 20:00. Training runs in 150 rounds of six minutes. Every device is at home at 05:00, and during the day about 35\% are, while the hub holds up to 32 devices. About 14\% of the devices are inside two cells at a time. S2 replaces the synthetic paths with 50 weekday user-days from the GeoLife trajectories~\cite{zheng2010geolife} and places five edges by $k$-means over the traces. S2 has a weaker commute pattern, and about 60\% of its devices stay at home during the day. S1-fast moves the S1 users at vehicle speed, and partial participation lets a random half of the devices train in each round. Two further settings change the class mix of all devices together. In the stepwise change, cell membership is fixed while the share parameter $\rho$ defined below rises from 0 to 0.8 and returns to 0 in five phases. In random mobility, devices move between cells by a Gauss--Markov model~\cite{liang1999gaussmarkov}, and $\rho$ changes from 0 to 0.8 halfway through the run. CIFAR-100 repeats S1 with 100 classes, and ResNet-18 repeats the stepwise change with a ResNet-18~\cite{he2016resnet} split in the middle. Section~\ref{sec:scale} repeats the S1 layout to reach 200 and 500 devices.

\begin{table}[t]
\centering
\caption{Evaluation settings. All settings except the scale study use 50 devices and five cells.}
\label{tab:settings}
\begin{adjustbox}{max width=\columnwidth}
\begin{tabular}{l>{\raggedright\arraybackslash}p{4.3cm}lc}
\toprule
Setting & Movement and class mix & Data, model & Seeds \\
\midrule
S1 (commute) & Walking commute between four residential cells and a hub & CIFAR-10, CNN & 5 \\
S2 (GeoLife) & 50 user-days of GeoLife traces & CIFAR-10, CNN & 3 \\
S1-fast & S1 at vehicle speed & CIFAR-10, CNN & 3 \\
Partial participation & S1, half of the devices train per round & CIFAR-10, CNN & 3 \\
Stepwise change & Fixed cells, $\rho$: 0, 0.4, 0.8, 0.4, 0 for all devices & CIFAR-10, CNN & 5 \\
Random mobility & Gauss--Markov movement, $\rho$: 0 then 0.8 & CIFAR-10, CNN & 3 \\
CIFAR-100 & S1 with 100 classes & CIFAR-100, CNN & 3 \\
ResNet-18 & Stepwise change & CIFAR-10, ResNet-18 & 3 \\
Scale & S1 layout repeated, 200 and 500 devices & CIFAR-10, CNN & 3 \\
\bottomrule
\end{tabular}
\end{adjustbox}
\end{table}

\textbf{Data and requests.} Each device holds two own classes of CIFAR-10~\cite{krizhevsky2009cifar}, or 20 of CIFAR-100. In S1, a device holds 620 to 2,857 training images. The requests of a device in an evaluation round are test images. For every Main request, a device receives on average $\rho$ OOP requests and $0.3\rho$ OOR requests. In the mobility settings, $\rho=0.1$ in the device's home cell and $\rho=0.8$ elsewhere, so OOP and OOR requests make up about 11.5\% and 51\% of the traffic.

\textbf{Models and training.} The CNN has a device block of 0.13\,M parameters, a device exit of 0.15\,M parameters, and an edge block and exit of 1.38\,M parameters. Its intermediate feature has $128\times8\times8$ values (32\,KB in single precision). All models are trained from random initialization with SGD (learning rate 0.01, weight decay $10^{-4}$, batch size 32, three local epochs per round) and the procedure of Section~\ref{sec:training} with $\lambda=0.4$ and $\Lambda=0.5$. Because training starts at the beginning of the day, the first evaluation rounds also cover the early phase of training.

\textbf{Inference rules.} All rules answer the same requests with the same trained models.
\begin{itemize}
\item \emph{SplitGP} answers with the device exit if its entropy is at most 0.8 nats and otherwise with the edge exit, as in confidence-based offloading~\cite{han2023splitgp}.
\item \emph{Device only} and \emph{Edge only} always use one exit.
\item \emph{Average} answers with the mean of the two probability vectors.
\item \emph{FedRoD} adds the logits of the two exits, which is the combination rule of FedRoD~\cite{chen2022fedrod}.
\item \emph{FedTHE} uses, for each request, the exit with the lower entropy. This is the solution of FedTHE's entropy objective over the ensemble weight~\cite{jiang2023fedthe}, because the entropy of a mixture is concave in the weight and is minimized at one end. FedTHE's feature-alignment term does not apply, because the two exits use different features.
\item \emph{ZTW} answers with the device exit when its largest probability exceeds a threshold and otherwise with the geometric mean of both exits~\cite{wolczyk2021ztw}. The threshold gives the same offloading ratio as SplitGP.
\item \emph{Label-shift EM} corrects the edge exit with the share of other classes estimated by EM on the device's recent requests and answers with the corrected edge exit~\cite{saerens2002adjusting}.
\item \emph{Learned weight} combines the exits with a per-request weight from a logistic regression on nine label-free features, such as the confidence of each exit and their agreement. The regression is trained with labels from separate development runs, which gives per-request weighting an advantage that a deployed system would not have.
\end{itemize}
DriftGate uses $r=1/2$ and a window of at least eight requests. In the simulation, the window holds the earlier requests of the current evaluation round, or the requests of the previous one when fewer than eight have arrived. Except in Section~\ref{sec:eval_offload}, DriftGate, Average, FedRoD, Label-shift EM, and Learned weight offload every request, and ZTW offloads the same share as SplitGP. The value of $r$, the window size, and the logistic regression were fixed on development runs whose seeds are not used in this section.

\textbf{Metric.} We report the accuracy of Eq.~\eqref{eq:accuracy} over the whole run. Evaluation takes place every five rounds in the mobility settings and every ten rounds in the stepwise change and random mobility. We report the mean over seeds, and differences are computed per seed before averaging.

\subsection{Overall Accuracy}
\label{sec:overall}

\begin{table*}[t]
\centering
\caption{Accuracy (\%) over the whole run. The last three rows give the gain of DriftGate over confidence-based offloading (SplitGP), its gain over the most accurate other rule in each setting, and that rule. Gains are the mean and standard deviation of the per-seed difference.}
\label{tab:main}
\begin{adjustbox}{max width=\textwidth}
\begin{tabular}{lcccccccc}
\toprule
Rule & S1 & S2 & S1-fast & Partial part. & Stepwise & Random mob. & CIFAR-100 & ResNet-18 \\
\midrule
SplitGP & 62.87 & 65.18 & 63.33 & 56.45 & 65.17 & 58.56 & 35.84 & 59.21 \\
Device only & 64.16 & 69.04 & 64.39 & 61.75 & 68.93 & 67.71 & 33.93 & 66.39 \\
Edge only & 62.91 & 65.24 & 63.38 & 56.44 & 65.20 & 58.61 & 35.84 & 57.04 \\
Average & 67.32 & 69.88 & 68.06 & 62.64 & 71.44 & 69.58 & 38.15 & 66.39 \\
FedRoD & 67.37 & 70.07 & 68.12 & 62.94 & 71.39 & 69.23 & 39.83 & 66.00 \\
FedTHE & 64.35 & 66.65 & 64.90 & 58.10 & 68.42 & 64.22 & 35.89 & 66.54 \\
ZTW & 67.36 & 70.06 & 68.11 & 62.94 & 71.38 & 69.22 & 39.83 & 65.99 \\
Label-shift EM & 64.84 & 68.63 & 65.41 & 60.45 & 67.59 & 65.42 & 40.44 & 61.73 \\
Learned weight & 66.81 & 68.59 & 67.52 & 62.02 & 70.69 & 68.55 & 37.54 & 63.61 \\
DriftGate & \textbf{68.31} & \textbf{71.70} & \textbf{68.96} & \textbf{64.82} & \textbf{72.36} & \textbf{71.09} & \textbf{41.29} & \textbf{66.64} \\
\midrule
Gain over SplitGP & 5.45 (0.47) & 6.52 (2.77) & 5.63 (0.38) & 8.37 (0.15) & 7.18 (0.77) & 12.53 (1.08) & 5.45 (0.67) & 7.43 (2.17) \\
Gain over best other & 0.95 (0.57) & 1.63 (1.04) & 0.84 (0.25) & 1.88 (0.59) & 0.92 (0.26) & 1.52 (0.53) & 0.85 (0.31) & 0.10 (0.61) \\
Best other rule & FedRoD & FedRoD & FedRoD & FedRoD & Average & Average & Label-shift EM & FedTHE \\
\bottomrule
\end{tabular}
\end{adjustbox}
\end{table*}

Is DriftGate more accurate than the existing ways of answering requests? Table~\ref{tab:main} compares the ten rules in the eight settings. The rules fall into two groups. Rules that choose one exit per request, namely SplitGP and FedTHE, do not rise clearly above the better single exit and often fall below it. Rules that combine the exits, namely Average, FedRoD, and ZTW, are 2 to 11 points above SplitGP. DriftGate is the most accurate rule in all eight settings. Its gain over SplitGP ranges from 5.4 points in S1 and CIFAR-100 to 12.5 points in random mobility, and its gain over the most accurate other rule ranges from 0.1 to 1.9 points.

The strongest alternative changes from setting to setting. It is FedRoD in the four settings built on individual movement, Average when the class mix changes for all devices together, Label-shift EM in CIFAR-100, and FedTHE with ResNet-18. No existing rule is the best choice across settings, whereas DriftGate is the most accurate in every one, with the largest margins under partial participation (1.9 points) and in S2 (1.6 points). In CIFAR-100, where each device owns 20 of 100 classes, the strongest alternative is Label-shift EM, which also corrects the edge exit. This confirms that the class bias of the edge exit matters, and DriftGate is 0.9 points above it because it combines the corrected edge exit with the device exit. The smallest gain appears with ResNet-18, whose large device block makes the device exit more accurate than the edge exit (66.4\% against 57.0\%). With little to gain from the edge, every reasonable rule ends close to the device exit, and DriftGate remains the most accurate by 0.1 points.

The combining rules also show why per-request weighting is not the answer. The learned weight uses labels from development runs to predict, for each request, which exit to trust. It is still 1.5 points below DriftGate in S1 and 3.1 points below in S2, because label-free features do not identify the kind of a single request reliably. A weight that reflects the device's recent traffic is enough, provided that the edge exit is corrected first.

\subsection{Where the Gain Comes From}
\label{sec:where}

\begin{figure*}[t]
\centering
\includegraphics[width=0.92\textwidth]{home_away}
\caption{Accuracy of devices in their home cell against accuracy of devices in another cell for six inference rules in S1 (left) and S2 (right), over the whole day.}
\label{fig:home_away}
\end{figure*}

\begin{table}[t]
\centering
\caption{Accuracy (\%) by request kind in S1 and S2 over the whole day.}
\label{tab:kinds}
\begin{adjustbox}{max width=\columnwidth}
\begin{tabular}{llccc}
\toprule
Setting & Rule & Main & OOP & OOR \\
\midrule
\multirow{6}{*}{S1} & SplitGP & 72.31 & 40.41 & 22.37 \\
 & Device only & 86.31 & 7.49 & 3.30 \\
 & Edge only & 72.27 & 40.79 & 22.52 \\
 & Average & 82.36 & 29.34 & 15.45 \\
 & FedRoD & 83.01 & 27.72 & 14.76 \\
 & DriftGate & 86.94 & 21.63 & 10.68 \\
\midrule
\multirow{6}{*}{S2} & SplitGP & 73.46 & 46.39 & 8.72 \\
 & Device only & 84.65 & 15.29 & 1.69 \\
 & Edge only & 73.42 & 46.92 & 8.82 \\
 & Average & 80.83 & 38.50 & 6.18 \\
 & FedRoD & 81.32 & 37.11 & 5.85 \\
 & DriftGate & 84.70 & 30.63 & 4.86 \\
\bottomrule
\end{tabular}
\end{adjustbox}
\end{table}

Which devices benefit from DriftGate? Fig.~\ref{fig:home_away} places each rule by the accuracy of devices at home (horizontal axis) and away (vertical axis). The single exits sit at opposite corners. The device exit is accurate at home and poor away, and the edge exit and SplitGP are the reverse. DriftGate moves to the right of every other rule. In S1, devices at home reach 79.8\%, which is 2.3 points above the device exit alone and 2.5 points above FedRoD. Devices away reach 53.2\%, within 1.6 points of the edge exit alone. In S2 the pattern is the same. Devices at home are 2.0 points above the device exit alone, and devices away are 0.4 points below the edge exit alone. DriftGate thus keeps the strength of the device exit at home while giving up little of the edge exit's accuracy away.

Table~\ref{tab:kinds} shows the same result by request kind. On Main requests, DriftGate is as accurate as the device exit alone (86.9\% against 86.3\% in S1), whereas Average and FedRoD lose 3 to 4 points there. The prior correction is what closes this gap, because it stops the edge exit from pulling Main requests toward the classes of other devices. On OOP and OOR requests, DriftGate stays well above the device exit, though below the edge exit. Main requests make up most of a device's traffic over the day, so this trade raises overall accuracy.

\begin{figure*}[t]
\centering
\includegraphics[width=0.92\textwidth]{time_of_day}
\caption{Accuracy over the day in S1 for SplitGP, the device exit alone, the edge exit alone, and DriftGate. The shaded background shows the share of devices in their home cell (right axis). Dashed lines separate the pre-commute, commute, daytime, return, and evening periods.}
\label{fig:time}
\end{figure*}

When during the day does the gain appear? Fig.~\ref{fig:time} follows the accuracy of four rules through the day in S1. Before the commute, almost every device is at home and DriftGate follows the device exit (72.9\% against 73.2\%). In these early hours the edge exit is still in its first rounds of training, and SplitGP, which sends most requests to it, stays at 55.3\%. During the day, when most devices are away, the device exit falls below the edge exit, and DriftGate stays above both exits from the commute until the evening. Its lead over the better single exit is 4.0 points during the commute, 1.7 points during the day, 1.5 points on the way back, and 3.4 points in the evening. The commute is when a device's traffic changes fastest, and it is also when combining the two exits helps most.

\subsection{Accuracy and Offloading}
\label{sec:eval_offload}

\begin{figure*}[t]
\centering
\includegraphics[width=0.92\textwidth]{offload_curves}
\caption{Accuracy against the offloading ratio when only requests on which the device exit is uncertain are sent to the edge. SplitGP and DriftGate offload requests whose device-exit entropy exceeds a threshold, and ZTW offloads requests whose largest device-exit probability is below a threshold. SplitGP answers the offloaded requests with the edge exit, ZTW with the geometric mean of both exits, and DriftGate with its combination. The dotted line marks the offloading ratio of SplitGP at its default threshold of 0.8 nats.}
\label{fig:offload}
\end{figure*}

How much accuracy does DriftGate keep when fewer requests can be offloaded? Fig.~\ref{fig:offload} varies the offloading threshold of each rule and plots accuracy against the share of requests sent to the edge. The threshold is on the device-exit entropy for SplitGP and DriftGate and on the largest device-exit probability for ZTW. All curves start from the device exit alone at an offloading ratio of zero. For SplitGP, sending more requests to the edge lowers accuracy in both scenarios, because the additional requests are more often Main requests, which the device exit answers better than the edge exit. ZTW and DriftGate improve as they offload more, since both use the device exit's output when they combine. DriftGate is above ZTW at every offloading ratio of 0.3 or more in all eight settings.

The curves also show how much offloading DriftGate needs. In S1, DriftGate exceeds every other rule in Table~\ref{tab:main}, most of which use the edge for every request, once it offloads 48\% of requests. In S2 it needs 28\%. Across the eight settings, this point lies between 11\% and 50\% in seven settings. CIFAR-100 needs 85\%, because with 100 classes the device exit is weak on its own and most requests benefit from the edge. SplitGP, by comparison, offloads 92\% of the requests in S1 and 84\% in S2 at its default threshold. In seven of the eight settings, DriftGate can therefore halve the traffic to the edge and still be more accurate than any rule that uses the edge for every request.

\subsection{Ablation}
\label{sec:ablation}

\begin{table}[t]
\centering
\caption{Accuracy (\%) of DriftGate and variants that remove one design choice.}
\label{tab:ablation}
\begin{adjustbox}{max width=\columnwidth}
\begin{tabular}{lccc}
\toprule
Variant & S1 & Random mob. & ResNet-18 \\
\midrule
DriftGate & \textbf{68.31} & \textbf{71.09} & 66.64 \\
Fixed weight $w=0.2$ & 68.16 & 70.51 & 64.02 \\
Fixed weight $w=0.5$ & 68.13 & 70.96 & \textbf{66.71} \\
No correction, $w=0.2$ & 64.75 & 63.78 & 61.25 \\
No correction, $w=0.5$ (Average) & 67.32 & 69.58 & 66.39 \\
Correction and product & 67.90 & 70.78 & 66.59 \\
\bottomrule
\end{tabular}
\end{adjustbox}
\end{table}

Which parts of DriftGate produce the gain? Table~\ref{tab:ablation} removes one choice at a time in S1, random mobility, and ResNet-18. Removing the prior correction costs the most. With $w=0.2$, accuracy drops by 3.4 points in S1 and by 6.7 points in random mobility. Without the correction, the edge exit pulls Main requests toward the classes of other devices, and giving the edge exit more weight makes this worse.

The weight matters differently for each model. A fixed weight of 0.2 is 2.6 points below DriftGate on ResNet-18, where the device exit is the stronger classifier. A fixed weight of 0.5 suits ResNet-18 but is 0.1 to 0.2 points below DriftGate on the CNN settings. The label-free weight of Eq.~\eqref{eq:weight} is at least as accurate as the better fixed weight in the two CNN settings and 0.07 points below it on ResNet-18, without being tuned for either model. Replacing the average with a product, which is how FedRoD and ZTW combine their outputs, lowers accuracy in all three settings even with the correction in place. This matches the argument of Section~\ref{sec:combine}. A product lets the device exit's near-zero probabilities on unlearned classes overrule the edge exit.

\subsection{Scale and the Least Accurate Devices}
\label{sec:scale}

\begin{table}[t]
\centering
\caption{Accuracy (\%) as the number of devices grows. The last column gives the offloading ratio at which DriftGate first matches the accuracy of SplitGP.}
\label{tab:scale}
\begin{adjustbox}{max width=\columnwidth}
\begin{tabular}{rcccccc}
\toprule
Devices & Cells & SplitGP & DriftGate & Gain & \multicolumn{2}{c}{Offloading ratio} \\
\cmidrule(lr){6-7}
 & & & & & SplitGP & DriftGate \\
\midrule
50 & 5 & 62.87 & 68.31 & 5.45 & 0.920 & 0.000 \\
200 & 20 & 65.26 & 69.00 & 3.73 & 0.949 & 0.218 \\
500 & 50 & 67.03 & 70.11 & 3.08 & 0.955 & 0.299 \\
\bottomrule
\end{tabular}
\end{adjustbox}
\end{table}

Does the gain hold in larger deployments? Table~\ref{tab:scale} repeats the S1 layout with 200 and 500 devices. DriftGate stays 3.1 to 5.4 points above SplitGP. With 200 and 500 devices, the device exit alone is less accurate than SplitGP, so DriftGate needs some offloading to match it. Even with 500 devices, DriftGate matches the accuracy of SplitGP while offloading 30\% of requests instead of 95.5\%. The edge work per request does not depend on the number of devices, because the correction uses only the class counts of the device's cells.

\begin{table}[t]
\centering
\caption{Accuracy (\%) of the 10\% least accurate devices, averaged over evaluation rounds.}
\label{tab:bottom}
\begin{adjustbox}{max width=\columnwidth}
\begin{tabular}{lcccc}
\toprule
Setting & SplitGP & Average & FedRoD & DriftGate \\
\midrule
S1 & 40.98 & 46.62 & 46.38 & \textbf{48.00} \\
S2 & 39.53 & 45.89 & 45.88 & \textbf{47.43} \\
S1-fast & 41.61 & 47.99 & 47.66 & \textbf{49.37} \\
Partial participation & 31.28 & 39.27 & 39.67 & \textbf{43.24} \\
\bottomrule
\end{tabular}
\end{adjustbox}
\end{table}

Does the average gain hide devices that do worse? Table~\ref{tab:bottom} reports the accuracy of the 10\% least accurate devices in each evaluation round. DriftGate raises this accuracy by 7.0 to 12.0 points over SplitGP and by 1.4 to 3.6 points over the strongest combining rule. The gain therefore reaches the devices that are served worst, not only the average device.

\subsection{Cost on Mobile Devices}
\label{sec:device}

\todo{Replace this subsection with the measured values. The structure and the expected values below follow the latency model of Section~\ref{sec:cost}; all red numbers are estimates.}

\textbf{Setup.} We measure the inference path on two smartphones, \todo{Samsung Galaxy model 1} and \todo{model 2}, and on an NVIDIA Jetson \todo{model} board. The device block and the device exit run with \todo{runtime, e.g., PyTorch Mobile or TensorRT}. The edge server is a desktop with an RTX 4090 GPU. We measure latency \todo{over $N$ requests after warm-up} and energy with \todo{power monitor or Android battery API}. The network is \todo{Wi-Fi / LTE}, with an uplink of \todo{measured} Mbps and a round-trip time of \todo{measured} ms.

\begin{table}[t]
\centering
\caption{Per-request cost of the inference path for the CNN. \textcolor{red}{Red values are estimates to be replaced by measurements.}}
\label{tab:device_cost}
\begin{adjustbox}{max width=\columnwidth}
\begin{tabular}{lcc}
\toprule
Component & Smartphone & Jetson \\
\midrule
Device block and device exit (ms) & \est{11.8} & \est{2.9} \\
Device block and device exit (mJ) & \est{TBD} & \est{TBD} \\
Uplink payload & 32\,KB & 32\,KB \\
Uplink time at 50\,Mbps (ms) & 5.2 & 5.2 \\
Edge block, exit, and correction (ms) & \est{0.6} & \est{0.6} \\
Downlink payload & 44\,B & 44\,B \\
Combination on the device ($\mu$s) & \est{5} & \est{3} \\
\bottomrule
\end{tabular}
\end{adjustbox}
\end{table}

\textbf{Per-request cost.} Table~\ref{tab:device_cost} lists the cost of each step. The device-side cost of DriftGate beyond SplitGP is the combination, which takes \est{a few microseconds}, and a downlink of 44 bytes instead of one label. Both are negligible next to the device block and the uplink of the intermediate feature.

\textbf{End-to-end latency and energy.} The mean latency of a request is the device-side time plus, for an offloaded request, the uplink time, the round-trip time, and the edge time. With the expected values of Table~\ref{tab:device_cost} and a round-trip time of 20\,ms, SplitGP in S1 takes \est{35.5\,ms} per request at an offloading ratio of 0.92 and reaches 62.9\% accuracy. DriftGate with every request offloaded takes \est{37.6\,ms} and reaches 68.3\%. With an offloading ratio of 0.5, DriftGate takes \est{24.7\,ms}, about \est{30\%} less than SplitGP, and reaches 67.5\%, which is above every other rule in Table~\ref{tab:main}. Energy follows the same pattern, because radio transmission dominates the energy of an offloaded request. \todo{Report measured energy per request for the three operating points and plot accuracy against latency and against energy.}
